% ============================================================
%  tick_invariance_note.tex
%  Why the (fixed) queue-lateness reward is only ASYMPTOTICALLY
%  tick-rate invariant, and not exactly invariant like the binary
%  reward: a Poisson-counting / Jensen's-inequality account.
% ============================================================
\documentclass[11pt,a4paper]{article}

\usepackage{amsmath,amssymb,amsthm}
\usepackage{geometry}
\geometry{margin=2.5cm}
\usepackage{booktabs}
\usepackage{array}
\usepackage{hyperref}
\usepackage{xcolor}
\usepackage{microtype}
\usepackage{parskip}

\newcommand{\E}{\mathbb{E}}
\newcommand{\Var}{\mathrm{Var}}
\newcommand{\Prob}{\mathbb{P}}

\newtheorem{remark}{Remark}

\title{Why queue-lateness is only \emph{asymptotically}\\tick-rate invariant}
\author{Technical note}
\date{}

\begin{document}
\maketitle

\begin{abstract}
The queue-lateness reward (\texttt{reward\_type}~1/3) was corrected with an
extra factor $1/\gamma$ so that, at fixed tick rate $\gamma$, the objective is
a positive constant rescaling of the original --- leaving every previously
reported ratio (FIFO/RVI, seed success) exactly unchanged.  What the fix does
\emph{not} achieve is exact invariance \emph{across} tick rates: measured FIFO
cost is $1497 / 1090 / 975 / 806 / 841 / 849$ at
$\gamma = 1/2/4/8/16/32$ --- converging, but not flat.  This note derives why,
precisely.  The mechanism is not vague ``discretisation noise'' but a specific,
quantifiable Jensen's-inequality effect arising because the tick counter is a
Poisson process sampling continuous physical time.  The same mechanism
explains the residual $O(1/\gamma)$ non-invariance already documented for the
tardiness-flux reward (\texttt{reward\_type}~4).
\end{abstract}

\tableofcontents

% ============================================================
\section{Setting}
% ============================================================

Write $\gamma$ for the tick rate, $D$ for a type's deadline in physical time
units, $d = D\gamma$ for the deadline expressed in ticks (an integer, in the
cases considered here).  Write $\tau$ for a job's physical waiting time and
$W$ for its recorded \emph{tick-count} waiting level --- the quantity the MDP
state actually tracks.

\begin{remark}[The tick clock is a single, state-independent Poisson process]
Inspection of the implementation confirms
\texttt{compute\_total\_tick\_rate} returns the constant \texttt{tick\_rate}
($=\gamma$), independent of how many queues are currently non-empty.  This
matches the exit-rate formula of the model, which has a single additive
$+\gamma$ term.  Consequently tick events occur as one global Poisson process
of rate $\gamma$, and every firing increments the recorded age of \emph{every}
currently non-empty queue by one tick.
\end{remark}

Because a given type-$n$ job's FIL clock advances at every global tick for as
long as that queue remains non-empty (i.e.\ for its entire waiting tenure),
the number of ticks recorded during a physical waiting duration $\tau$ is
exactly the number of Poisson($\gamma$) events observed in an interval of
length $\tau$:
\begin{equation}
  W \mid \tau \;\sim\; \mathrm{Poisson}(\gamma\tau)
  \qquad\text{\emph{exactly}, not as an approximation.}
  \label{eq:poisson}
\end{equation}

\begin{remark}[FIFO's physical trajectory does not depend on $\gamma$]
Under FIFO, a job enters service exactly when a compatible server becomes
free while it is the oldest waiting job of a servable type --- a decision
triggered only by arrivals and service completions.  Tick events change
recorded ages but never create or remove a serving opportunity for FIFO.  The
physical sample path (who is served, and when) is therefore identical for
every $\gamma$; the tick rate affects only how finely that physical path is
\emph{measured} into an integer count.  This is what licenses treating $\tau$
in \eqref{eq:poisson} as a $\gamma$-independent random variable with a fixed
distribution (namely, FIFO's true physical excess-waiting-time distribution),
while $W$'s conditional law depends on $\gamma$ only through the counting
relationship above.
\end{remark}

% ============================================================
\section{The mean is unbiased; the noise is not free}
% ============================================================

From \eqref{eq:poisson},
\begin{equation}
  \E[W \mid \tau] = \gamma\tau,
  \qquad
  \Var(W \mid \tau) = \gamma\tau,
  \label{eq:moments}
\end{equation}
using that a Poisson distribution has variance equal to its mean.  The
coefficient of variation is
\begin{equation}
  \mathrm{CV}(W\mid\tau) \;=\; \frac{\sqrt{\gamma\tau}}{\gamma\tau}
  \;=\; \frac{1}{\sqrt{\gamma\tau}} \;\xrightarrow[\gamma\to\infty]{} \;0.
  \label{eq:cv}
\end{equation}
At $\gamma = 1$ this ratio is large: $W$ is a \emph{noisy} proxy for the
continuous quantity $\gamma\tau$.  At $\gamma = 32$ the same physical duration
is measured far more precisely, in relative terms.

The cost function does not use $W$ linearly.  It applies
$B = \max(0, W - d)$ --- a \textbf{convex} function of $W$ --- for the FIL
term, and a further convex (roughly quadratic) transform for the tail term.
By Jensen's inequality, for any convex $g$,
\begin{equation}
  \E[g(W)] \;\geq\; g(\E[W]).
  \label{eq:jensen}
\end{equation}
So $\E[B]$ is \textbf{systematically inflated} above the naively expected
value $\max(0,\gamma\tau - d) = \gamma\max(0,\tau - D)$.  The size of this gap
is governed by the variance in \eqref{eq:moments}, and it is exactly why
measured FIFO cost decreases monotonically toward its asymptotic value as
$\gamma$ grows: the Poisson noise that inflates the convex transform shrinks
in relative terms.

% ============================================================
\section{Quantifying the gap, term by term}
% ============================================================

\subsection{The linear (FIL) term}

Write $\tau_{\text{excess}} = \max(0, \tau - D)$ for the true physical excess.
For $W$ far above the threshold $d$ (deep in the tail, i.e.\ substantially
late), $B \approx W - d$ and the Jensen gap is small in relative terms.  Near
the threshold --- i.e.\ for jobs close to their deadline --- approximating the
large-mean Poisson by a Normal distribution with
$\mu = \gamma\tau$, $\sigma = \sqrt{\gamma\tau}$, the standard expected-shortfall
(``loss function'') identity gives
\begin{equation}
  \E[\max(0, W - d)]
  \;\approx\;
  \underbrace{\gamma\,\tau_{\text{excess}}}_{\text{leading term, } O(\gamma)}
  \;+\;
  \underbrace{\sigma\,\phi(z)}_{\text{boundary correction, } O(\sqrt{\gamma})},
  \qquad \sigma = \sqrt{\gamma\tau},
  \label{eq:loss-fn}
\end{equation}
where $\phi$ is the standard normal density and $z$ measures how many standard
deviations the mean sits from the threshold.  The boundary correction is
$O(\sqrt\gamma)$: negligible \emph{relative to} the leading $O(\gamma)$ term
once $\gamma$ is large, but not negligible in absolute size at small $\gamma$
--- and it contributes most from states whose physical wait is close to $D$,
which is exactly the population of states that matters most for a
deadline-sensitive policy.

\subsection{The quadratic (tail) term}

This correction has a cleaner closed form.  Using
$\Var(B) \approx \Var(W) = \gamma\tau$ for states clearly past the deadline
(where the thresholding no longer distorts the variance materially),
\begin{equation}
  \E[B^2] \;=\; \Var(B) + \E[B]^2
  \;\approx\; \gamma\tau \;+\; \gamma^2\tau_{\text{excess}}^2,
\end{equation}
so that
\begin{equation}
  \boxed{\;\frac{\E[B^2]}{\gamma^2}
  \;\approx\; \tau_{\text{excess}}^2 \;+\; \frac{\tau}{\gamma}\;}
  \qquad\xrightarrow[\gamma\to\infty]{}\;\tau_{\text{excess}}^2.
  \label{eq:quad}
\end{equation}
The residual term $\tau/\gamma$ is exactly the Poisson
variance-to-mean-squared ratio from \eqref{eq:cv}, carried through the
quadratic transform.  This is the precise, closed-form source of the
$O(1/\gamma)$ correction in the tail (Koole) term of the queue-lateness
reward.

% ============================================================
\section{Why $\gamma = 1$ shows a much larger gap than $\gamma = 4$}
% ============================================================

Take $D = 3$ (our benchmark deadline).  At $\gamma = 1$, $d = 3$ ticks: the
deadline is reached after only three counting events, so
$\mathrm{CV}(W \mid \tau{=}D) = 1/\sqrt{d} = 1/\sqrt{3} \approx 0.58$ --- the
tick count is a coarse, high-relative-noise measurement of physical time right
where it matters most.  At $\gamma = 4$, $d = 12$ ticks, and
$\mathrm{CV} = 1/\sqrt{12} \approx 0.29$: four times the counting resolution
for the same physical deadline, hence roughly half the relative noise, and
correspondingly a much smaller Jensen gap.  This matches the direction and
approximate scale of the measured drop from $1497$ ($\gamma=1$) to $975$
($\gamma=4$).

% ============================================================
\section{Summary statement}
% ============================================================

\begin{equation}
  g^{\text{QL}}_{\text{measured}}(\gamma)
  \;=\;
  g^{\text{QL}}_{\text{continuous}}
  \;+\; O\!\left(\gamma^{-1/2}\right) \text{ or } O\!\left(\gamma^{-1}\right)
  \text{-type corrections},
  \qquad
  \xrightarrow[\gamma\to\infty]{} \; 0.
\end{equation}

The mechanism is entirely a \emph{counting-noise and convexity} story.  The
tick rate does not change the physical system at all --- FIFO's true dynamics
are $\gamma$-independent, as established above --- it only changes how
precisely a Poisson tick-count can measure continuous elapsed time.  Because
the cost function applies convex (threshold, square) transforms to that noisy
count, Jensen's inequality leaks a $\gamma$-dependent bias back into the
otherwise-unbiased mean.  This is structurally the \emph{same} mechanism
already documented for the tardiness-flux reward (\texttt{reward\_type}~4),
whose residual non-invariance comes from applying a crossing indicator to the
same discretely-sampled continuous quantity.

\begin{remark}[Honest caveat on the fit]
The numbers above ($1497/1090/975/806/841/849$ at
$\gamma=1/2/4/8/16/32$) come from a single Monte Carlo evaluation
($2\times10^5$ periods) at each tick rate --- enough to confirm the direction
and general shape of convergence, not enough to distinguish a
dominant $O(\gamma^{-1})$ correction from $O(\gamma^{-1/2})$ in this regime.
Pinning that down would need a longer evaluation horizon and a denser
log-spaced sweep of $\gamma$.
\end{remark}

\end{document}
